% GENERATED FILE. Edit canonical lessons, then run npm run generate:books.
% canonical-source-hash: fnv1a64:7a96945adf9fb4d2
% canonical-lessons: UR-C273-se-bara, UR-C273-behtar, UR-C273-subah, UR-W273-suad, UR-C273-do-sal-se, UR-R273-comparing-and-time

\chapter{Comparing, and How Long}
\label{ch:ur-comparing-and-how-long}
\hlchaptermodality{273}

\begin{chapteropening}
\textbf{By the end of this chapter:} \emph{I can compare two things (bas ṭaiksī se sastī hai, sab se achchhā), and say how long or since when (do sāl se, subah se shām tak, din bhar).}

Complete UR-R273-comparing-and-time: compare two things, and say how long.
\end{chapteropening}

\section[\texorpdfstring{\ur{سے} \ur{بڑا} (se baṛā)}{سے بڑا (se baṛā)}]{\ur{سے} \ur{بڑا} (se baṛā) — bigger than, the biggest}

\label{lesson:UR-C273-se-bara}



\begin{quote}
Say \textquotedblleft{}big\textquotedblright{} and \textquotedblleft{}small\textquotedblright{}. (\emph{baṛā}, \emph{chhoṭā}.) Then \textquotedblleft{}I liked this city\textquotedblright{}. (\emph{mujhe yih shahr achchhā lagā}.) Now compare two cities.
\end{quote}

\subsection*{You'll want to know: bigger than}
\begin{quote}
\textbf{\ur{لاہور} \ur{بڑا} \ur{ہے۔}}
\end{quote}

\begin{quote}
\emph{Lāhaur baṛā hai.} — \textquotedblleft{}Lahore is big.\textquotedblright{}
\end{quote}

\begin{quote}
\textbf{\ur{یہ} \ur{گھر} \ur{اس} \ur{گھر} \ur{سے} \ur{بڑا} \ur{ہے۔}}
\end{quote}

\begin{quote}
\emph{yih ghar us ghar se baṛā hai.} — \textquotedblleft{}This house is bigger than that one.\textquotedblright{}
\end{quote}

Urdu adds no \textquotedblleft{}-er\textquotedblright{} here. \textbf{\ur{سے}} (\emph{se}), \textquotedblleft{}from, than\textquotedblright{}, goes after the thing compared, and the word stays as it is:

\begin{quote}
\textbf{\ur{بس} \ur{ٹیکسی} \ur{سے} \ur{سستی} \ur{ہے۔}}
\end{quote}

\begin{quote}
\emph{bas ṭaiksī se sastī hai.} — \textquotedblleft{}The bus is cheaper than a taxi.\textquotedblright{}
\end{quote}

\textbf{\ur{بس}} is feminine: \textbf{\ur{سستی}}. \textbf{\ur{زیادہ}} (\emph{zyādā}) and \textbf{\ur{کم}} (\emph{kam}) say more and less:

\begin{quote}
\textbf{\ur{مجھے} \ur{چائے} \ur{زیادہ} \ur{پسند} \ur{ہے۔}}
\end{quote}

\begin{quote}
\emph{mujhe chāy zyādā pasand hai.} — \textquotedblleft{}I like tea more.\textquotedblright{}
\end{quote}

\subsection*{You'll want to know: the most}
\begin{quote}
\textbf{\ur{یہ} \ur{سب} \ur{سے} \ur{اچھا} \ur{ہے۔}}
\end{quote}

\begin{quote}
\emph{yih sab se achchhā hai.} — \textquotedblleft{}This is the best.\textquotedblright{}
\end{quote}

\begin{quote}
\textbf{\ur{وہ} \ur{سب} \ur{سے} \ur{لمبی} \ur{لڑکی} \ur{ہے۔}}
\end{quote}

\begin{quote}
\emph{voh sab se lambī laṛkī hai.} — \textquotedblleft{}She is the tallest girl.\textquotedblright{}
\end{quote}

\textbf{\ur{سب} \ur{سے}} (\emph{sab se}), \textquotedblleft{}than all\textquotedblright{}, makes the most: \textbf{\ur{سب} \ur{سے} \ur{بڑا}}, \textbf{\ur{سب} \ur{سے} \ur{سستا}}.

\subsection*{Your turn}
Say these aloud:

\begin{itemize}
\raggedright
  \item \textquotedblleft{}This is the biggest\textquotedblright{} — \emph{yih sab se baṛā hai}
  \item \textquotedblleft{}This is the smallest\textquotedblright{} — \emph{yih sab se chhoṭā hai}
  \item \textquotedblleft{}This is the cheapest\textquotedblright{} — \emph{yih sab se sastā hai}
\end{itemize}

\begin{tcolorbox}[breakable,colback=teal!4,colframe=teal!35!black,title={Before you move on}]
\textquotedblleft{}Cheaper than a taxi\textquotedblright{}? (\textbf{\ur{ٹیکسی} \ur{سے} \ur{سستی}}.) \textquotedblleft{}The best\textquotedblright{}? (\textbf{\ur{سب} \ur{سے} \ur{اچھا}}.) Where does \textbf{\ur{سے}} go? (After the thing compared.)
\end{tcolorbox}
\section[\texorpdfstring{\ur{بہتر} (behtar)}{بہتر (behtar)}]{\ur{بہتر} (behtar) — better}

\label{lesson:UR-C273-behtar}



\begin{quote}
Before the new one: say \textquotedblleft{}the bus is cheaper than a taxi\textquotedblright{}. (\emph{bas ṭaiksī se sastī hai}.) \textquotedblleft{}Good\textquotedblright{} has a word of its own for \textquotedblleft{}better\textquotedblright{}.
\end{quote}

\subsection*{You'll want to know: \ur{بہتر}}
\textbf{\ur{بہتر}} — \emph{behtar} — \textquotedblleft{}better\textquotedblright{}.

\begin{quote}
\textbf{\ur{یہ} \ur{بہتر} \ur{ہے۔}}
\end{quote}

\begin{quote}
\emph{yih behtar hai.} — \textquotedblleft{}This is better.\textquotedblright{}
\end{quote}

\begin{quote}
\textbf{\ur{بس} \ur{ٹیکسی} \ur{سے} \ur{بہتر} \ur{ہے۔}}
\end{quote}

\begin{quote}
\emph{bas ṭaiksī se behtar hai.} — \textquotedblleft{}The bus is better than a taxi.\textquotedblright{}
\end{quote}

\begin{quote}
\textbf{\ur{آج} \ur{میں} \ur{بہتر} \ur{ہوں۔}}
\end{quote}

\begin{quote}
\emph{āj maiṅ behtar hūṅ.} — \textquotedblleft{}I'm better today.\textquotedblright{}
\end{quote}

It does not change for a man or a woman, and it takes \textbf{\ur{سے}} like any comparison.

\begin{cousinweb}
\textbf{\ur{بہتر}} (\emph{behtar}) is usually traced to Persian \emph{behtar}, \textquotedblleft{}better\textquotedblright{}, built from \emph{beh}, \textquotedblleft{}good\textquotedblright{}, and the Persian ending \emph{-tar}, \textquotedblleft{}more\textquotedblright{}. Urdu took it whole.
\end{cousinweb}

\subsection*{Your turn}
Say these aloud:

\begin{itemize}
\raggedright
  \item \emph{behtar}, then \emph{yih behtar hai}
  \item \textquotedblleft{}Tea is better than coffee\textquotedblright{} — \emph{chāy kāfī se behtar hai}
\end{itemize}

\begin{tcolorbox}[breakable,colback=teal!4,colframe=teal!35!black,title={Before you move on}]
What does \textbf{\ur{بہتر}} mean? (\textquotedblleft{}Better\textquotedblright{}.) \textquotedblleft{}I'm better today\textquotedblright{}? (\textbf{\ur{آج} \ur{میں} \ur{بہتر} \ur{ہوں}}.)
\end{tcolorbox}
\section[\texorpdfstring{\ur{صبح} (subah)}{صبح (subah)}]{\ur{صبح} (subah) — the morning}

\label{lesson:UR-C273-subah}



\begin{quote}
Before the new one: say \textquotedblleft{}this is better\textquotedblright{}. (\emph{yih behtar hai}.) Then \textquotedblleft{}evening\textquotedblright{}. (\emph{shām}.) This lesson gives the other end of the day.
\end{quote}

\subsection*{You'll want to know: \ur{صبح}}
\textbf{\ur{صبح}} — \emph{subah} — \textquotedblleft{}the morning\textquotedblright{}. It is feminine.

\begin{quote}
\textbf{\ur{صبح} \ur{بارش} \ur{ہوئی۔}}
\end{quote}

\begin{quote}
\emph{subah bārish huī.} — \textquotedblleft{}It rained in the morning.\textquotedblright{}
\end{quote}

\begin{quote}
\textbf{\ur{میں} \ur{صبح} \ur{جلدی} \ur{اٹھتا} \ur{ہوں۔}}
\end{quote}

\begin{quote}
\emph{maiṅ subah jaldī uṭhtā hūṅ.} — \textquotedblleft{}I get up early in the morning.\textquotedblright{}
\end{quote}

Like \textbf{\ur{شام}} (\emph{shām}), it can stand alone as a time: \textquotedblleft{}in the morning\textquotedblright{}.

\begin{cousinweb}
\textbf{\ur{صبح}} (\emph{subah}) is usually traced, through Persian, to Arabic \emph{ṣubḥ}, \textquotedblleft{}morning, dawn\textquotedblright{}. Its first letter is new to your hand; the next lesson takes it out.
\end{cousinweb}

\subsection*{Your turn}
Say these aloud:

\begin{itemize}
\raggedright
  \item \emph{subah}, then \emph{subah bārish huī}
  \item \textquotedblleft{}I get up early in the morning\textquotedblright{} — \emph{maiṅ subah jaldī uṭhtā hūṅ}
\end{itemize}

\begin{tcolorbox}[breakable,colback=teal!4,colframe=teal!35!black,title={Before you move on}]
What does \textbf{\ur{صبح}} mean? (\textquotedblleft{}The morning\textquotedblright{}.) Its partner at the end of the day? (\textbf{\ur{شام}}.)
\end{tcolorbox}
\section[\texorpdfstring{\ur{ص} (suād)}{ص (suād)}]{\ur{ص} (suād) — the letter that opens \ur{صبح}}

\label{lesson:UR-W273-suad}



\begin{quote}
Before the new one: write \ur{ڈ} and \ur{ض} once each, and name them. (\emph{ḍāl}, \emph{zuād}.)

Now say \textbf{\ur{صبح}} (\emph{subah}), \textquotedblleft{}the morning\textquotedblright{}. This lesson takes one letter out of it and puts it in your hand.
\end{quote}

\subsection*{Script you'll notice: \ur{ص}}
\textbf{\ur{ص}} — \emph{suād}, the sound \emph{s}, the same as \textbf{\ur{س}}.

It opens \textbf{\ur{صبح}} (\emph{subah}). You already write \textbf{\ur{ض}} (\emph{zuād}): \textbf{\ur{ص}} is the same body \textbf{without its dot}, an oval loop with a short tooth. At the start of a word, as here, the loop joins the next letter; standing alone, it ends in a bowl below the line. Urdu keeps it for words from Arabic.

\subsection*{Writing: \ur{ص}}
Put your pen on \ur{ص} and follow its line. Copy the shape you can see — slowly, and larger than it is printed.

\begin{quote}
This book does not yet tell you \textbf{where to start this letter or which way to travel}. It is not written down here until it can be written down with a source. Copying what is in front of you needs no such source, and it is how the shape gets into your hand in the meantime.
\end{quote}

\subsection*{Your turn}
\begin{itemize}
\raggedright
  \item \emph{Look:} at \textbf{\ur{صبح}} and point at \ur{ص} inside it
  \item \emph{Trace it:} \ur{ص} three times, saying \emph{s} as you finish each one
  \item \emph{Write it:} \textbf{\ur{صبح}} once, joined, with no dot on its first letter
  \item \emph{Read it:} \textbf{\ur{صبح}} aloud once more
\end{itemize}

\begin{tcolorbox}[breakable,colback=teal!4,colframe=teal!35!black,title={Before you move on}]
Which letter is this — \ur{ص}? (\textbf{suād}, \emph{s}.) What tells it from \ur{ض}? (No dot.) Which word did it come from? (\textbf{\ur{صبح}}, \emph{subah}, \textquotedblleft{}the morning\textquotedblright{}.)
\end{tcolorbox}
\section[\texorpdfstring{\ur{دو} \ur{سال} \ur{سے} (do sāl se)}{دو سال سے (do sāl se)}]{\ur{دو} \ur{سال} \ur{سے} (do sāl se) — how long, and since when}

\label{lesson:UR-C273-do-sal-se}



\begin{quote}
Say \textquotedblleft{}the morning\textquotedblright{} and \textquotedblleft{}a year\textquotedblright{}. (\emph{subah}, \emph{sāl}.) Then \textquotedblleft{}two\textquotedblright{}. (\emph{do}.)
\end{quote}

\subsection*{You'll want to know: since, and for}
\begin{quote}
\textbf{\ur{دو} \ur{سال} \ur{سے۔}}
\end{quote}

\begin{quote}
\emph{do sāl se.} — \textquotedblleft{}For two years.\textquotedblright{}
\end{quote}

\begin{quote}
\textbf{\ur{میں} \ur{دو} \ur{سال} \ur{سے} \ur{یہاں} \ur{رہتا} \ur{ہوں۔}}
\end{quote}

\begin{quote}
\emph{maiṅ do sāl se yahāṅ rahtā hūṅ.} — \textquotedblleft{}I've lived here for two years.\textquotedblright{}
\end{quote}

\textbf{\ur{سے}} (\emph{se}), \textquotedblleft{}from\textquotedblright{}, after a time says since when, or for how long. Urdu keeps the present: the living began two years ago and goes on.

\begin{quote}
\textbf{\ur{وہ} \ur{کل} \ur{سے} \ur{بیمار} \ur{ہے۔}}
\end{quote}

\begin{quote}
\emph{voh kal se bīmār hai.} — \textquotedblleft{}She's been ill since yesterday.\textquotedblright{}
\end{quote}

\subsection*{You'll want to know: from … to, all day}
\begin{quote}
\textbf{\ur{میں} \ur{صبح} \ur{سے} \ur{شام} \ur{تک} \ur{کام} \ur{کرتا} \ur{ہوں۔}}
\end{quote}

\begin{quote}
\emph{maiṅ subah se shām tak kām kartā hūṅ.} — \textquotedblleft{}I work from morning to evening.\textquotedblright{}
\end{quote}

\textbf{\ur{سے}} opens the span and \textbf{\ur{تک}} (\emph{tak}), \textquotedblleft{}up to\textquotedblright{}, closes it. \textbf{\ur{بھر}} (\emph{bhar}), \textquotedblleft{}full\textquotedblright{}, after a time word fills it:

\begin{quote}
\textbf{\ur{کل} \ur{دن} \ur{بھر} \ur{بارش} \ur{ہوئی۔}}
\end{quote}

\begin{quote}
\emph{kal din bhar bārish huī.} — \textquotedblleft{}Yesterday it rained all day.\textquotedblright{}
\end{quote}

\textbf{\ur{رات} \ur{بھر}} (\emph{rāt bhar}) is \textquotedblleft{}all night\textquotedblright{}.

\subsection*{Your turn}
Say these aloud:

\begin{itemize}
\raggedright
  \item \textquotedblleft{}I've lived here for two years\textquotedblright{} — \emph{maiṅ do sāl se yahāṅ rahtā hūṅ}
  \item \textquotedblleft{}I've lived here for three years\textquotedblright{} — \emph{maiṅ tīn sāl se yahāṅ rahtā hūṅ}
  \item \textquotedblleft{}I've lived here for a year\textquotedblright{} — \emph{maiṅ ek sāl se yahāṅ rahtā hūṅ}
\end{itemize}

\begin{tcolorbox}[breakable,colback=teal!4,colframe=teal!35!black,title={Before you move on}]
\textquotedblleft{}Since yesterday\textquotedblright{}? (\textbf{\ur{کل} \ur{سے}}.) \textquotedblleft{}From morning to evening\textquotedblright{}? (\textbf{\ur{صبح} \ur{سے} \ur{شام} \ur{تک}}.) \textquotedblleft{}All day\textquotedblright{}? (\textbf{\ur{دن} \ur{بھر}}.)
\end{tcolorbox}
\section[(review)]{Review — comparing, and how long}

\label{lesson:UR-R273-comparing-and-time}



\begin{quote}
Say \textquotedblleft{}I want to become a teacher\textquotedblright{} and \textquotedblleft{}I'm happy\textquotedblright{}. (\emph{maiṅ ustād bannā chāhtā hūṅ}, \emph{maiṅ khush hūṅ}.)
\end{quote}

\subsection*{Your turn}
\begin{itemize}
\raggedright
  \item \emph{Say it:} The bus is cheaper than a taxi. In my view, it is better.
  \item \emph{Say it:} This is the best tea.
  \item \emph{Say it:} I've lived in Lahore for two years.
  \item \emph{Say it:} I work from morning to evening.
  \item \emph{Say it:} Yesterday it rained all day.
  \item \emph{Write it:} \textbf{\ur{صبح}} once, and point at \ur{ص} inside it
\end{itemize}

\begin{tcolorbox}[breakable,colback=teal!4,colframe=teal!35!black,title={Before you move on}]
\textquotedblleft{}Bigger than this\textquotedblright{}? (\textbf{\ur{اس} \ur{سے} \ur{بڑا}}.) Now the other way round: what does \textbf{\ur{سب} \ur{سے}} add? (The most.) \textquotedblleft{}Since yesterday\textquotedblright{}? (\textbf{\ur{کل} \ur{سے}}.)
\end{tcolorbox}
